% ============================================================
% MÓDULO: EL RECIBO DE SUELDO Y LA LIQUIDACIÓN SALARIAL
% ============================================================
\clearpage
\section*{MÓDULO: EL RECIBO DE SUELDO Y LA LIQUIDACIÓN SALARIAL}
\subsection*{Remuneración, beneficios y negociación}

\begin{recuadro}
\textbf{Objetivo del módulo:} saber leer y controlar un recibo de sueldo, entender de dónde
sale cada número, conocer los beneficios de un trabajador registrado (aguinaldo, vacaciones,
obra social, ART) y cómo se negocian los salarios.
\end{recuadro}

% ------------------------------------------------------------
\section*{1. EL TRABAJO REGISTRADO}

Un trabajador \textbf{registrado} (``en blanco'') tiene un contrato de trabajo regido por la
\textbf{Ley de Contrato de Trabajo (LCT, Ley 20.744)}, modificada en 2026 por la
\textbf{Ley 27.802 de Modernización Laboral}, y por el \textbf{convenio colectivo de trabajo
(CCT)} de su actividad. Está dado de alta en ARCA, tiene aportes jubilatorios, obra social,
cobertura de la ART y cobra con \textbf{recibo de sueldo}.

\paragraph*{Sueldo bruto y sueldo neto.}
\begin{itemize}
  \item \textbf{Bruto:} el total de lo que se gana antes de los descuentos.
  \item \textbf{Neto:} lo que efectivamente se cobra (bruto $-$ descuentos).
\end{itemize}
Cuando se negocia un sueldo en una entrevista, conviene aclarar siempre si se habla de
\textbf{bruto o neto}.

% ------------------------------------------------------------
\section*{2. LOS CONCEPTOS DEL RECIBO}

\subsubsection*{2.1. Haberes remunerativos}
Son parte del salario: generan aportes y contribuciones y se tienen en cuenta para el
aguinaldo, las vacaciones, la indemnización y la jubilación.
\begin{itemize}
  \item \textbf{Sueldo básico:} el fijado por el convenio para la categoría (por ejemplo,
        ``técnico'', ``oficial especializado'').
  \item \textbf{Antigüedad:} un adicional por cada año trabajado, según el convenio.
  \item \textbf{Presentismo:} premio por asistencia y puntualidad, si el convenio lo prevé.
  \item \textbf{Horas extra:} se pagan con un recargo del 50\% en días hábiles y del 100\% en
        sábados después de las 13 h, domingos y feriados.
  \item \textbf{Adicionales:} por título, por tarea riesgosa, por turno nocturno, comisiones.
\end{itemize}

\subsubsection*{2.2. Conceptos no remunerativos}
No generan aportes ni (en general) se usan para calcular aguinaldo y vacaciones: viáticos con
comprobante, algunas sumas acordadas en paritarias como ``no remunerativas'', asignaciones
familiares.

\subsubsection*{2.3. Deducciones: los aportes del trabajador}
\begin{center}
\small
\begin{tabular}{|p{0.45\linewidth}|c|p{0.30\linewidth}|}
\hline
\textbf{Aporte} & \textbf{\%} & \textbf{Destino} \\ \hline
Jubilación (SIPA) & 11\% & Tu futura jubilación \\ \hline
Ley 19.032 (INSSJP -- PAMI) & 3\% & Obra social de los jubilados \\ \hline
Obra social & 3\% & Tu cobertura de salud \\ \hline
Cuota sindical & según convenio & Sindicato de la actividad \\ \hline
Impuesto a las Ganancias & si corresponde & Retención para ARCA (sueldos altos) \\ \hline
\end{tabular}
\end{center}

\subsubsection*{2.4. Contribuciones patronales}
Las paga el empleador \textbf{además} del sueldo bruto: jubilación, PAMI, obra social, Fondo
Nacional de Empleo, asignaciones familiares, ART y seguro de vida colectivo obligatorio. Desde
2026, con la reglamentación de la Ley 27.802 (Decreto 407/2026), el nuevo recibo debe mostrar
también estas contribuciones y el \textbf{costo laboral total} que paga la empresa por cada
trabajador.

% ------------------------------------------------------------
\section*{3. EJEMPLO DE LIQUIDACIÓN}

\textit{Técnico con 3 años de antigüedad; básico \$1.200.000; antigüedad 1\% por año;
presentismo \$100.000; 10 horas extra al 50\% (valor hora = básico / 200 h).}

\begin{center}
\small
\begin{tabular}{|p{0.52\linewidth}|r|r|}
\hline
\textbf{Concepto} & \textbf{Haberes} & \textbf{Deducciones} \\ \hline
Sueldo básico & \$1.200.000 & \\ \hline
Antigüedad (3 años $\times$ 1\%) & \$36.000 & \\ \hline
Presentismo & \$100.000 & \\ \hline
Horas extra 50\% (10 h $\times$ \$9.000) & \$90.000 & \\ \hline
\textbf{Total remunerativo (bruto)} & \textbf{\$1.426.000} & \\ \hline
Jubilación 11\% & & \$156.860 \\ \hline
Ley 19.032 (PAMI) 3\% & & \$42.780 \\ \hline
Obra social 3\% & & \$42.780 \\ \hline
Cuota sindical 2\% & & \$28.520 \\ \hline
\textbf{Total deducciones} & & \textbf{\$270.940} \\ \hline
\multicolumn{2}{|r|}{\textbf{SUELDO NETO A COBRAR}} & \textbf{\$1.155.060} \\ \hline
\end{tabular}
\end{center}

\paragraph*{Qué controlar en el recibo.}
Datos personales y CUIL; fecha de ingreso (de ella depende la antigüedad); categoría y
convenio; período liquidado; que los porcentajes sean correctos; fecha y lugar de pago. El
recibo se firma (o se acepta digitalmente) como constancia de que se cobró, no de que se está
de acuerdo: si hay un error, se puede firmar ``en disconformidad'' y reclamar.

% ------------------------------------------------------------
\section*{4. BENEFICIOS DEL TRABAJADOR REGISTRADO}

\subsubsection*{4.1. Aguinaldo (Sueldo Anual Complementario, SAC)}
Equivale al \textbf{50\% de la mejor remuneración mensual} de cada semestre. Se paga en dos
cuotas: la primera hasta el 30 de junio y la segunda hasta el 18 de diciembre. Si no se
trabajó el semestre completo, se paga en proporción a los días trabajados.

\textit{Ejemplo:} si el mejor sueldo bruto del semestre fue \$1.426.000, el aguinaldo bruto es
\$713.000 (y se le descuentan los mismos aportes).

\subsubsection*{4.2. Vacaciones}
\begin{center}
\small
\begin{tabular}{|p{0.45\linewidth}|c|}
\hline
\textbf{Antigüedad} & \textbf{Días corridos} \\ \hline
Menos de 5 años & 14 \\ \hline
De 5 a 10 años & 21 \\ \hline
De 10 a 20 años & 28 \\ \hline
Más de 20 años & 35 \\ \hline
\end{tabular}
\end{center}
Para tener derecho al período completo hay que haber trabajado al menos la mitad de los días
hábiles del año; si no, corresponde un día de vacaciones cada 20 trabajados. Las vacaciones se
pagan por adelantado (el día de vacaciones vale el sueldo mensual dividido 25). Con la Ley
27.802 pueden fraccionarse en períodos de al menos una semana, por acuerdo entre las partes.

\subsubsection*{4.3. Obra social, ART y licencias}
\begin{itemize}
  \item \textbf{Obra social:} cobertura médica para el trabajador y su familia; se puede
        elegir (opción de cambio de obra social).
  \item \textbf{ART (Aseguradora de Riesgos del Trabajo):} cubre los accidentes de trabajo
        (en el lugar de trabajo o \textit{in itinere}, en el trayecto entre la casa y el
        trabajo) y las enfermedades profesionales. Ante un accidente: avisar al empleador y
        pedir la denuncia a la ART.
  \item \textbf{Licencias:} por enfermedad, por nacimiento, por matrimonio, por examen, por
        fallecimiento de un familiar, según la ley y el convenio.
\end{itemize}

% ------------------------------------------------------------
\section*{5. NEGOCIACIÓN SALARIAL}

\subsubsection*{5.1. Negociación colectiva: las paritarias}
Los sindicatos y las cámaras empresarias de cada actividad negocian los
\textbf{convenios colectivos} y sus actualizaciones salariales (\textbf{paritarias}). El
acuerdo fija los salarios mínimos de cada categoría; una empresa puede pagar más, pero no
menos. Existe además el \textbf{Salario Mínimo, Vital y Móvil}, piso para todos los
trabajadores.

\subsubsection*{5.2. Negociación individual: cómo pedir un sueldo}
\begin{itemize}
  \item \textbf{Informarse:} conocer el convenio de la actividad y los sueldos de mercado para
        el puesto (portales de empleo, encuestas salariales, colegas).
  \item \textbf{Hablar en bruto mensual} y preguntar qué incluye: convenio, adicionales,
        obra social, horario, horas extra, capacitación, viáticos.
  \item \textbf{Dar un rango razonable} y justificarlo con lo que aportás (licencias,
        experiencia, idiomas, manejo de software).
  \item \textbf{Mirar el paquete completo:} un sueldo algo menor con capacitación paga,
        habilitaciones o crecimiento puede convenir más al inicio de la carrera.
  \item \textbf{Para un aumento:} pedir una reunión, llevar logros concretos y medibles y
        proponer, no exigir.
\end{itemize}

% ------------------------------------------------------------
\section*{6. ACTIVIDADES DEL MÓDULO}

\begin{enumerate}
  \item \textbf{Leer un recibo real:} con un recibo de sueldo (con datos tapados) aportado por
        un familiar o el docente, identificar cada concepto y verificar los porcentajes.
  \item \textbf{Liquidar un sueldo:} en una planilla de cálculo, liquidar el sueldo de un
        técnico con antigüedad, horas extra y presentismo, el aguinaldo y las vacaciones.
  \item \textbf{Bruto vs. neto:} si te ofrecen \$1.500.000 brutos, ¿cuánto cobrás? ¿Y cuánto le
        cuesta a la empresa aproximadamente?
  \item \textbf{Convenio:} buscar el convenio colectivo de una actividad donde podría
        trabajar un técnico aviónico (talleres, aerolíneas, metalúrgica) y la escala salarial
        vigente.
  \item \textbf{Simulacro de negociación:} en parejas, uno hace de empleador y otro de
        candidato que negocia su primer sueldo.
\end{enumerate}
